\documentclass[11pt,a4paper]{article}
\usepackage[utf8]{inputenc}
\usepackage[margin=1in]{geometry}
\usepackage{amsmath,amssymb,amsfonts}
\usepackage{booktabs}
\usepackage{graphicx}
\usepackage{hyperref}
\usepackage{microtype}
\usepackage{siunitx}
\usepackage{cite}
\usepackage{tikz}
\usepackage{xcolor}

\hypersetup{
    colorlinks=true,
    linkcolor=blue!70!black,
    citecolor=blue!70!black,
    urlcolor=blue!70!black
}

\title{\textbf{WindTunnel lab: Physical and Mathematical Foundations of the 2D Lattice Boltzmann Wind Tunnel}}
\author{\textbf{North Carolina School of Science and Mathematics (NCSSM)} \\ High-Performance Computational Aerodynamics Laboratory}
\date{\today}

\begin{document}
\maketitle

\begin{abstract}
This document provides a comprehensive mathematical and physical formulation of \textbf{WindTunnel lab}, a real-time, two-dimensional numerical wind tunnel developed at NCSSM based on the Lattice Boltzmann Method (LBM) with a D2Q9 lattice topology. The document details the kinetic Boltzmann equation, the Bhatnagar-Gross-Krook (BGK) collision operator, Chapman-Enskog multiscale asymptotic analysis recovering the incompressible Navier-Stokes equations, Smagorinsky subgrid-scale (LES) turbulence modeling, aerodynamic boundary conditions, momentum exchange force evaluation for lift ($C_L$) and drag ($C_D$) coefficients, and the MacCormack predictor-corrector transport scheme for high-contrast, anti-diffusive smoke streamline visualization.
\end{abstract}

\tableofcontents
\newpage

\section{Introduction and Governing Equations}

In aerodynamic wind tunnel testing, external flow past an immersed lifting body is governed by the incompressible Navier-Stokes equations expressing conservation of mass and linear momentum:
\begin{align}
\nabla \cdot \mathbf{u} &= 0, \label{eq:continuity} \\
\frac{\partial \mathbf{u}}{\partial t} + (\mathbf{u} \cdot \nabla)\mathbf{u} &= -\frac{1}{\rho_0} \nabla p + \nu \nabla^2 \mathbf{u}, \label{eq:navier_stokes}
\end{align}
where $\mathbf{u} = (u_x, u_y)^T$ denotes the macroscopic velocity vector, $p$ is the static pressure field, $\rho_0$ is the reference fluid density, and $\nu$ is the kinematic viscosity.

The fundamental dimensionless parameter governing aerodynamic similarity is the Reynolds number based on the airfoil chord length $c$ and freestream velocity $U_\infty$:
\begin{equation}
\mathrm{Re} = \frac{U_\infty c}{\nu}.
\end{equation}

Rather than discretizing Equations~\eqref{eq:continuity}--\eqref{eq:navier_stokes} via traditional continuum methods (finite differences, finite volumes, or finite elements), WindTunnel lab employs the \textbf{Lattice Boltzmann Method} (LBM). LBM simulates fluid mechanics from a mesoscopic kinetic perspective, solving for the evolution of discrete particle distribution functions whose statistical velocity moments recover the Navier-Stokes equations in the low-Mach, nearly-incompressible limit.

\section{The Lattice Boltzmann Method (LBM D2Q9)}

\subsection{Discrete Boltzmann Equation}
The core governing equation for the discrete velocity distribution functions $f_i(\mathbf{x}, t)$ along discrete lattice direction $\mathbf{c}_i$ is the discrete Boltzmann equation with the Bhatnagar-Gross-Krook (BGK) single relaxation time approximation:
\begin{equation}
f_i(\mathbf{x} + \mathbf{c}_i \Delta t, t + \Delta t) - f_i(\mathbf{x}, t) = -\frac{\Delta t}{\tau} \left[ f_i(\mathbf{x}, t) - f_i^{(0)}(\mathbf{x}, t) \right], \quad i = 0, \dots, 8,
\label{eq:lbm_bgk}
\end{equation}
where $f_i^{(0)}(\mathbf{x}, t)$ is the local Maxwell-Boltzmann equilibrium distribution function, $\Delta t$ is the discrete time step, and $\tau$ is the dimensionless relaxation time parameter.

Equation~\eqref{eq:lbm_bgk} is typically split into two distinct, highly parallelizable operations:
\begin{enumerate}
    \item \textbf{Collision}:
    \begin{equation}
    f_i^*(\mathbf{x}, t) = f_i(\mathbf{x}, t) - \frac{1}{\tau} \left[ f_i(\mathbf{x}, t) - f_i^{(0)}(\mathbf{x}, t) \right].
    \end{equation}
    \item \textbf{Streaming}:
    \begin{equation}
    f_i(\mathbf{x} + \mathbf{c}_i \Delta t, t + \Delta t) = f_i^*(\mathbf{x}, t).
    \end{equation}
\end{enumerate}

\subsection{D2Q9 Velocity Topology and Weights}
The two-dimensional nine-velocity (D2Q9) lattice model comprises one rest particle, four orthogonal nearest-neighbor particles, and four diagonal next-nearest-neighbor particles:
\begin{equation}
\mathbf{c}_i = \begin{cases}
(0, 0), & i = 0, \\
\left(\cos \frac{(i-1)\pi}{2}, \, \sin \frac{(i-1)\pi}{2}\right) c, & i = 1, 2, 3, 4, \\
\sqrt{2}\left(\cos \left[\frac{(i-5)\pi}{2} + \frac{\pi}{4}\right], \, \sin \left[\frac{(i-5)\pi}{2} + \frac{\pi}{4}\right]\right) c, & i = 5, 6, 7, 8,
\end{cases}
\end{equation}
where $c = \Delta x / \Delta t = 1$ in standard lattice units. Specifically:
\begin{equation}
\mathbf{C} = \begin{bmatrix}
0 & 1 & 0 & -1 & 0 & 1 & -1 & -1 & 1 \\
0 & 0 & 1 & 0 & -1 & 1 & 1 & -1 & -1
\end{bmatrix}^T.
\end{equation}

The corresponding lattice weights $w_i$, derived from Hermite polynomial quadrature of the continuous Maxwellian, are:
\begin{equation}
w_i = \begin{cases}
\frac{4}{9}, & i = 0, \\[4pt]
\frac{1}{9}, & i = 1, 2, 3, 4, \\[4pt]
\frac{1}{36}, & i = 5, 6, 7, 8.
\end{cases}
\end{equation}

The lattice speed of sound $c_s$ satisfies:
\begin{equation}
c_s^2 = \frac{1}{3} c^2 = \frac{1}{3} \frac{\Delta x^2}{\Delta t^2}.
\end{equation}

\subsection{Equilibrium Distribution Function}
A second-order Taylor-Hermite series expansion of the Maxwellian velocity distribution around zero mean velocity yields the local equilibrium distribution:
\begin{equation}
f_i^{(0)}(\rho, \mathbf{u}) = w_i \rho \left[ 1 + \frac{\mathbf{c}_i \cdot \mathbf{u}}{c_s^2} + \frac{(\mathbf{c}_i \cdot \mathbf{u})^2}{2 c_s^4} - \frac{\mathbf{u} \cdot \mathbf{u}}{2 c_s^2} \right].
\end{equation}
Substituting $c_s^2 = 1/3$:
\begin{equation}
f_i^{(0)}(\rho, \mathbf{u}) = w_i \rho \left[ 1 + 3 (\mathbf{c}_i \cdot \mathbf{u}) + \frac{9}{2} (\mathbf{c}_i \cdot \mathbf{u})^2 - \frac{3}{2} (\mathbf{u} \cdot \mathbf{u}) \right].
\end{equation}

\subsection{Macroscopic Moments and Equation of State}
The hydrodynamic conserved macroscopic quantities are the local density $\rho$ and momentum density $\rho \mathbf{u}$:
\begin{align}
\rho(\mathbf{x}, t) &= \sum_{i=0}^8 f_i(\mathbf{x}, t), \\
\rho \mathbf{u}(\mathbf{x}, t) &= \sum_{i=0}^8 \mathbf{c}_i f_i(\mathbf{x}, t).
\end{align}

The effective thermodynamic pressure follows an ideal gas equation of state:
\begin{equation}
p = c_s^2 \rho = \frac{1}{3} \rho.
\end{equation}

\subsection{Chapman-Enskog Asymptotic Expansion}
Using a multi-scale asymptotic expansion with small Knudsen number $\epsilon = \Delta t \ll 1$,
\begin{align}
f_i &= f_i^{(0)} + \epsilon f_i^{(1)} + \epsilon^2 f_i^{(2)} + \mathcal{O}(\epsilon^3), \\
\partial_t &= \epsilon \partial_{t_1} + \epsilon^2 \partial_{t_2}, \quad \nabla = \epsilon \nabla_1,
\end{align}
one systematically recovers the continuity and Navier-Stokes equations at $\mathcal{O}(\epsilon^2)$, with kinematic viscosity given by:
\begin{equation}
\nu = c_s^2 \left( \tau - \frac{1}{2} \right) \Delta t = \frac{2\tau - 1}{6} \Delta t.
\end{equation}
For a specified physical Reynolds number $\mathrm{Re}$, chord length $c$, and lattice inlet velocity $U_\infty = u_{\text{in}}$, the relaxation time is computed dynamically as:
\begin{equation}
\nu = \frac{u_{\text{in}} c}{\mathrm{Re}}, \quad \tau = 0.5 + 3\nu.
\end{equation}
Positivity of physical viscosity requires $\tau > 0.5$.

\section{Turbulence Modeling: Smagorinsky Subgrid-Scale (LES) Model}

At elevated Reynolds numbers, numerical stability of single-relaxation-time LBM can degrade when $\tau \to 0.5$. To ensure robust simulation of high-speed wake turbulence and separated shear layers, WindTunnel lab incorporates the \textbf{Smagorinsky Large Eddy Simulation (LES)} subgrid-scale model.

\subsection{Resolved Strain Rate Tensor}
In standard continuum LES, unresolved turbulent eddies are represented by an effective eddy viscosity $\nu_t$:
\begin{equation}
\nu_t = (C_s \Delta)^2 |\bar{S}|,
\end{equation}
where $C_s \approx 0.15$ is the Smagorinsky constant, $\Delta = \Delta x = 1$ is the grid filter width, and $|\bar{S}| = \sqrt{2 \bar{S}_{\alpha\beta} \bar{S}_{\alpha\beta}}$ is the characteristic filtered strain rate tensor:
\begin{equation}
\bar{S}_{\alpha\beta} = \frac{1}{2} \left( \frac{\partial u_\alpha}{\partial x_\beta} + \frac{\partial u_\beta}{\partial x_\alpha} \right).
\end{equation}

\subsection{Non-Equilibrium Stress Tensor Formulation}
A profound advantage of LBM is that the strain rate tensor does \textbf{not} require spatial finite differences; it is evaluated purely locally from the second-order non-equilibrium moment of the distribution functions:
\begin{equation}
\Pi_{\alpha\beta}^{(1)} = \sum_{i=0}^8 c_{i\alpha} c_{i\beta} \left( f_i - f_i^{(0)} \right).
\end{equation}
From Chapman-Enskog expansion, the leading-order non-equilibrium tensor is directly proportional to the strain rate:
\begin{equation}
\Pi_{\alpha\beta}^{(1)} = -2 \rho c_s^2 \tau_{\text{eff}} \bar{S}_{\alpha\beta}.
\end{equation}
Taking the Frobenius norm:
\begin{equation}
Q = \sqrt{2 \Pi_{\alpha\beta}^{(1)} \Pi_{\alpha\beta}^{(1)}} = \sqrt{\Pi_{xx}^2 + \Pi_{yy}^2 + 2\Pi_{xy}^2} = 2 \rho c_s^2 \tau_{\text{eff}} |\bar{S}|.
\end{equation}

\subsection{Effective Relaxation Time Calculation}
The total effective viscosity is the sum of molecular and subgrid eddy viscosities:
\begin{equation}
\nu_{\text{eff}} = \nu_0 + \nu_t = c_s^2 \left(\tau_{\text{eff}} - \frac{1}{2}\right).
\end{equation}
Substituting $\nu_t = C_s^2 |\bar{S}| = C_s^2 \frac{Q}{2 \rho c_s^2 \tau_{\text{eff}}}$ leads to a closed quadratic equation for $\tau_{\text{eff}}$:
\begin{equation}
\tau_{\text{eff}} = \tau_0 + \frac{C_s^2 Q}{2 \rho c_s^4 \tau_{\text{eff}}}.
\end{equation}
Solving the positive root of this quadratic equation yields:
\begin{equation}
\tau_{\text{eff}} = \frac{1}{2} \left[ \tau_0 + \sqrt{\tau_0^2 + 18 \sqrt{2} \, C_s^2 \frac{Q}{\rho}} \right].
\end{equation}
This local adaptation ensures that wherever turbulent shearing $Q$ increases, $\tau_{\text{eff}}$ increases safely away from the $0.5$ singularity, guaranteeing unconditional numerical stability.

\section{Boundary Conditions}

\subsection{Inflow Boundary Condition (Inlet at $x=0$)}
At the wind tunnel inlet ($x = 0$), uniform horizontal freestream velocity $\mathbf{u}_{\text{in}} = (U_\infty, 0)$ is prescribed. An equilibrium Dirichlet boundary condition is applied across all incoming directions ($i \in \{1, 5, 8\}$):
\begin{equation}
f_i(0, y, t) = f_i^{(0)}(\rho=1.0, u_x=U_\infty, u_y=0).
\end{equation}

\subsection{Outflow Boundary Condition (Outlet at $x=L_x$)}
At the tunnel outlet ($x = L_x$), a convective zero-gradient boundary condition allows vortices and pressure waves to exit without reflection back into the test section:
\begin{equation}
f_i(L_x, y, t) = f_i(L_x - 1, y, t), \quad \forall i.
\end{equation}

\subsection{Tunnel Sidewalls ($y=0$ and $y=L_y$)}
To model an unconfined wind tunnel test section without artificial boundary layer growth, top and bottom walls employ free-slip symmetry conditions:
\begin{align}
f_{\text{UP}}(x, 0, t) &= f_{\text{DOWN}}(x, 0, t), \quad \text{at bottom wall } y = 0, \\
f_{\text{DOWN}}(x, L_y, t) &= f_{\text{UP}}(x, L_y, t), \quad \text{at top wall } y = L_y,
\end{align}
where mirror pairs exchange populations: $(2 \leftrightarrow 4)$, $(5 \leftrightarrow 8)$, $(6 \leftrightarrow 7)$.

\subsection{Airfoil No-Slip Wall: Half-Way Bounce-Back}
On the solid boundaries of the airfoil $\Omega_s$, the no-slip condition ($\mathbf{u} = 0$) is enforced via the half-way bounce-back scheme. When a streaming particle encounters a solid node, its momentum is exactly reversed into the conjugate direction:
\begin{equation}
f_{\bar{i}}(\mathbf{x}_b, t + \Delta t) = f_i^*(\mathbf{x}_b, t),
\end{equation}
where $\mathbf{c}_{\bar{i}} = -\mathbf{c}_i$ corresponds to the opposite velocity index mapping:
\begin{equation}
\mathrm{OPP} = [0, 3, 4, 1, 2, 7, 8, 5, 6].
\end{equation}

\section{Aerodynamic Force and Coefficient Evaluation}

\subsection{Momentum Exchange Method (MEM)}
The aerodynamic forces exerted on the immersed airfoil are calculated directly using the \textbf{Momentum Exchange Method}. At every solid node $\mathbf{x}_b \in \Omega_s$, fluid packets transfer net linear momentum during bounce-back:
\begin{equation}
\Delta \mathbf{p}(\mathbf{x}_b, t) = \mathbf{c}_i f_i^*(\mathbf{x}_b, t) - \mathbf{c}_{\bar{i}} f_{\bar{i}}(\mathbf{x}_b, t+\Delta t) = 2 \mathbf{c}_i f_i^*(\mathbf{x}_b, t).
\end{equation}
Summing over all solid nodes and incoming directions yields the total hydrodynamic force vector:
\begin{equation}
\mathbf{F}(t) = \sum_{\mathbf{x}_b \in \Omega_s} \sum_{i=1}^8 2 \mathbf{c}_i f_i^*(\mathbf{x}_b, t).
\end{equation}
In wind-tunnel axes aligned with the $+x$ freestream direction:
\begin{align}
\text{Total Drag Force: } & D = F_x, \\
\text{Total Lift Force: } & L = F_y.
\end{align}

\subsection{Dimensionless Aerodynamic Coefficients}
The dynamic pressure per unit length is:
\begin{equation}
q_\infty = \frac{1}{2} \rho_0 U_\infty^2 c.
\end{equation}
The instantaneous aerodynamic coefficients are:
\begin{equation}
C_L = \frac{L}{q_\infty} = \frac{F_y}{\frac{1}{2} \rho_0 U_\infty^2 c}, \quad
C_D = \frac{D}{q_\infty} = \frac{F_x}{\frac{1}{2} \rho_0 U_\infty^2 c}, \quad
C_R = \sqrt{C_L^2 + C_D^2}.
\end{equation}

\subsection{Surface Pressure Coefficient ($C_p$)}
The local dimensionless surface pressure coefficient $C_p$ is derived from the local density variation $\rho - \rho_0$:
\begin{equation}
C_p = \frac{p - p_\infty}{\frac{1}{2} \rho_0 U_\infty^2} = \frac{\frac{1}{3}(\rho - 1)}{\frac{1}{2} U_\infty^2}.
\end{equation}
At the stagnation point where $\mathbf{u} = 0$, Bernoulli's equation dictates $C_p = 1.0$.

\subsection{Unsteady Vortex Shedding, Cycle Averaging, and Stall Characterization}
Beyond the linear pre-stall regime ($\alpha > \alpha_{\text{stall}} \approx 13^\circ - 16^\circ$), boundary layer separation occurs over the suction surface, giving rise to an unsteady periodic K\'{a}rm\'{a}n vortex shedding street in the wake. Under these separated conditions, the instantaneous lift coefficient exhibits periodic or quasi-periodic fluctuations about a mean:
\begin{equation}
C_L(t) = \overline{C_L} + \Delta C_L \sin(\omega_s t + \psi) + \sum_{k=2}^\infty A_k \sin(k \omega_s t + \psi_k),
\label{eq:cl_fluctuation}
\end{equation}
where $\Delta C_L$ is the vortex shedding fluctuation amplitude, $\omega_s = 2\pi f_s$ is the shedding angular frequency, and $\psi$ is an arbitrary initial phase. The fundamental shedding frequency $f_s$ scales with the projected frontal height $h_{\text{proj}} = c \sin\alpha$ and freestream velocity $U_\infty$ according to the Strouhal number:
\begin{equation}
\mathrm{St} = \frac{f_s c \sin\alpha}{U_\infty} \approx 0.18 - 0.20.
\end{equation}

When student data acquisition relies on a single instantaneous snapshot sample $C_L(t_0)$, the measured value depends strongly on the phase $\psi(t_0) \in [0, 2\pi)$ at the instant the operator triggers the acquisition. Because post-stall shedding amplitudes can reach $\Delta C_L \approx \pm 0.3 - 0.5$, raw instantaneous samples exhibit severe scatter that destroys the coherence of the aerodynamic polar ($C_L \text{ vs. } \alpha$) and obscures the stall peak.

To eliminate phase variance, WindTunnel lab implements a \textbf{Cycle-Averaged Data Acquisition (DAQ) Burst}. When triggered, the system samples a 5.0-second burst ($N \approx 200$ simulation frames at 40 FPS), applies a 3-point filter to suppress high-frequency lattice noise, detects local extrema (peaks and troughs), and integrates over integer shedding cycles:
\begin{equation}
\overline{C_L} = \frac{1}{t_K - t_1} \int_{t_1}^{t_K} C_L(t) \, dt \approx \frac{1}{K - 1} \sum_{k=1}^{K-1} \left( \frac{1}{i_{k+1} - i_k} \sum_{j=i_k}^{i_{k+1}-1} C_{L, j} \right),
\end{equation}
where $t_1$ and $t_K$ correspond to the first and last detected shedding peaks. Because the integration domain spans exact integer cycles, the periodic fluctuating harmonic vanishes identically:
\begin{equation}
\int_{t_1}^{t_K} \sin(\omega_s t + \psi) \, dt \equiv 0.
\end{equation}

In addition to the cycle-mean $\overline{C_L}$ and $\overline{C_D}$, the acquisition protocol records the physical oscillation envelope $[C_{L,\min}, C_{L,\max}]$, fluctuation amplitude $\Delta C_L = (C_{L,\max} - C_{L,\min})/2$, and standard deviation $\sigma_{C_L}$. In the user interface and polar plot, this is visualized as vertical error whiskers and translucent confidence bands, preserving both the clean mean stall curve and the quantitative magnitude of wake turbulence.

\section{Smoke Streamline Transport: MacCormack Scheme}

To visualize streamlines and shedding wake streaklines with high fidelity, a conserved passive scalar field $\phi(\mathbf{x}, t) \in [0, 1]$ (smoke dye concentration) is advected by the local flow field:
\begin{equation}
\frac{\partial \phi}{\partial t} + \mathbf{u} \cdot \nabla \phi = 0.
\end{equation}

\subsection{The MacCormack Predictor-Corrector Anti-Diffusion Scheme}
Standard first-order semi-Lagrangian advection suffers from significant numerical diffusion:
\begin{equation}
\phi_{\text{diff}}(\mathbf{x}) \approx \phi(\mathbf{x} - \mathbf{u}\Delta t) = \phi(\mathbf{x}) - \Delta t (\mathbf{u}\cdot\nabla)\phi + \frac{\Delta t^2}{2} (\mathbf{u}\cdot\nabla)^2 \phi + \mathcal{O}(\Delta x),
\end{equation}
which turns sharp streamlines into a milky, washed-out fog.

WindTunnel lab replaces standard upwinding with a second-order \textbf{MacCormack predictor-corrector transport scheme}:
\begin{enumerate}
    \item \textbf{Forward Predictor}: Advect the dye forward along the velocity field:
    \begin{equation}
    \phi^* = \mathcal{S}(\phi^n, \, \mathbf{u}\Delta t),
    \end{equation}
    where $\mathcal{S}(\phi, \mathbf{d})$ denotes bilinear interpolation sampling at $\mathbf{x} - \mathbf{d}$.
    \item \textbf{Backward Reverser}: Advect the predicted field backward along $-\mathbf{u}\Delta t$:
    \begin{equation}
    \phi^{**} = \mathcal{S}(\phi^*, \, -\mathbf{u}\Delta t).
    \end{equation}
    \item \textbf{Error Estimation and Compensation}:
    The numerical error accumulated over the forward-backward cycle is $\mathcal{E} = \frac{1}{2}(\phi^n - \phi^{**})$. We construct the compensated intermediate field:
    \begin{equation}
    \phi^c = \operatorname{clip}\left(\phi^n + \frac{1}{2}(\phi^n - \phi^{**}), \, 0.0, \, 1.0\right).
    \end{equation}
    \item \textbf{Final Forward Advection}:
    \begin{equation}
    \phi^{n+1} = \mathcal{S}(\phi^c, \, \mathbf{u}\Delta t) \times \gamma_{\text{decay}},
    \end{equation}
    where $\gamma_{\text{decay}} = 0.9995$ is a subtle background dissipation factor that purges residual accumulation while preserving razor-sharp cores.
\end{enumerate}

\subsection{Gaussian Injector Rake and Interactive Wand}
At the inlet ($x = 0$), discrete dye injectors with Gaussian profiles prevent pixelation artifacts:
\begin{equation}
\phi_{\text{inlet}}(y) = \sum_{k=1}^{N_{\text{lines}}} \exp\left( -\frac{(y - y_k)^2}{2\sigma^2} \right), \quad \sigma = \max(1.1, \, 0.12 \Delta y_{\text{lines}}).
\end{equation}
In pulsed mode, injection is periodically modulated:
\begin{equation}
\phi_{\text{inject}}(t) = \begin{cases}
1, & (n_{\text{step}} \bmod 24) < 11, \\
0, & \text{otherwise.}
\end{cases}
\end{equation}
For the interactive smoke wand positioned at user coordinates $(x_w, y_w)$:
\begin{equation}
\phi_{\text{wand}}(x, y) = 0.9 \exp\left( -\frac{(x - x_w)^2 + (y - y_w)^2}{2 r_w^2} \right).
\end{equation}

\section{NACA 4-Digit Airfoil Geometry}

The geometry of the immersed lifting surface is parameterized by the standard NACA 4-digit designation ($MPXX$):
\begin{itemize}
    \item $M$: Maximum camber ($m = M / 100$).
    \item $P$: Location of maximum camber ($p = P / 10$).
    \item $XX$: Maximum thickness as percentage of chord ($t = XX / 100$).
\end{itemize}

\subsection{Thickness Distribution}
The symmetrical half-thickness distribution $y_t(x)$ along the chord coordinate $x/c \in [0, 1]$ is:
\begin{equation}
y_t(x) = 5 t c \left[ 0.2969 \sqrt{\frac{x}{c}} - 0.1260 \left(\frac{x}{c}\right) - 0.3516 \left(\frac{x}{c}\right)^2 + 0.2843 \left(\frac{x}{c}\right)^3 - 0.1036 \left(\frac{x}{c}\right)^4 \right].
\end{equation}

\subsection{Mean Camber Line}
For cambered profiles ($m > 0, p > 0$):
\begin{equation}
y_c(x) = \begin{cases}
\frac{m}{p^2} \left[ 2p\left(\frac{x}{c}\right) - \left(\frac{x}{c}\right)^2 \right] c, & 0 \le \frac{x}{c} < p, \\[8pt]
\frac{m}{(1-p)^2} \left[ (1-2p) + 2p\left(\frac{x}{c}\right) - \left(\frac{x}{c}\right)^2 \right] c, & p \le \frac{x}{c} \le 1.
\end{cases}
\end{equation}

\subsection{Quarter-Chord Pivot Rotation}
To accurately model the aerodynamic center and angle of attack $\alpha$, the airfoil coordinates are rotated around its quarter-chord pivot point:
\begin{equation}
\mathbf{x}_q = \left(x_{\text{le}} + 0.25c, \, y_{\text{le}}\right).
\end{equation}
Any lattice node $(X, Y)$ is transformed into the un-rotated airfoil frame $(x_r, y_r)$ via:
\begin{align}
x_r &= (X - x_q) \cos\alpha - (Y - y_q) \sin\alpha, \\
y_r &= (X - x_q) \sin\alpha + (Y - y_q) \cos\alpha.
\end{align}
A node is flagged as solid ($\Omega_s = \text{True}$) if:
\begin{equation}
0 \le x_r + 0.25c \le c \quad \text{and} \quad |y_r - y_c(x_r)| \le y_t(x_r).
\end{equation}

\section{Summary of Parameters and Units}

\begin{table}[htbp]
\centering
\caption{Key physical and numerical parameters in WindTunnel lab.}
\begin{tabular}{@{}llll@{}}
\toprule
\textbf{Symbol} & \textbf{Physical Meaning} & \textbf{Typical Value} & \textbf{Units} \\ \midrule
$N_x \times N_y$ & Grid resolution (medium) & $800 \times 300$ & Lattice nodes \\
$c$ & Airfoil chord length & $0.16 \times N_x \approx 128$ & Lattice spacing $\Delta x$ \\
$\alpha$ & Angle of attack (AoA) & $-20^\circ \text{ to } +35^\circ$ & Degrees \\
$U_\infty$ & Lattice inlet velocity & $0.08$ & Lattice velocity $\Delta x / \Delta t$ \\
$\mathrm{Re}$ & Reynolds number & $1000 - 30000$ & Dimensionless \\
$\tau$ & BGK relaxation time & $> 0.501$ & Lattice time steps $\Delta t$ \\
$c_s$ & Speed of sound in D2Q9 & $1/\sqrt{3} \approx 0.577$ & Lattice velocity \\
$C_s$ & Smagorinsky constant & $0.15$ & Dimensionless \\
$C_L, C_D$ & Sectional lift and drag & Computed dynamically & Dimensionless \\
\bottomrule
\end{tabular}
\end{table}

\section{Cloud Deployment and Google Colab Architecture}

To democratize computational aerodynamics education and eliminate the barrier of local GPU hardware requirements, WindTunnel lab includes a turnkey cloud integration targeting \textbf{Google Colab}.

\subsection{Architectural Challenge: Real-Time WebSockets}
In standard web deployments, reverse proxies route standard HTTP request-response cycles. However, WindTunnel lab streams 16-bit scalar lattice fields and dynamic force vectors at 30--40~FPS over a bidirectional WebSocket (\texttt{/ws}). Google Colab's native container proxy (\texttt{google.colab.kernel.proxyPort}) is restricted to standard HTTP/REST traffic and fails to negotiate the \texttt{101 Switching Protocols} handshake necessary for continuous binary WebSockets.

\subsection{Cloudflare Tunnel Integration}
To provide students with a zero-configuration, secure, high-bandwidth communication pipeline, the deployment architecture couples the Uvicorn/FastAPI backend with a transient \textbf{Cloudflare Tunnel} (\texttt{cloudflared}). When the Colab notebook executes:
\begin{enumerate}
    \item The LBM solver initializes in an asynchronous worker thread utilizing the allocated NVIDIA GPU (e.g., Tesla T4 or L4) via CuPy.
    \item The FastAPI server binds locally to \texttt{http://127.0.0.1:8000}.
    \item A background \texttt{cloudflared} agent opens a multiplexed secure tunnel to Cloudflare edge nodes, routing an ephemeral public HTTPS/WSS URL (e.g., \texttt{https://*.trycloudflare.com}) directly to the local server.
\end{enumerate}
Because the frontend automatically synchronizes with \texttt{window.location.protocol} (\texttt{https:} $\to$ \texttt{wss:}), the WebGL client establishes a persistent low-latency binary stream. This enables students on any device (chromebooks, laptops, tablets) to independently execute laboratory measurements, map lift-drag polars, and export CSV stall datasets directly to their local machines without any local software installation.

\section{Conclusion}
The combination of D2Q9 Lattice Boltzmann kinetics, Smagorinsky subgrid-scale turbulence closure, momentum exchange boundary force integration, MacCormack anti-diffusive advection, and cloud-native Colab deployment provides WindTunnel lab with a rigorous, physically validated foundation capable of reproducing complex aerodynamic phenomena---such as boundary layer separation, stall vortex shedding, and dynamic force responses---in real time at high lattice throughput.

\end{document}
